%% kapa: exam=1 round=27 year=2016 period=1 form=A subject=civil_law subject_ko=민법 seq=18 number=18
%% source: Q-Net 감정평가사 기출문제 (cst00309) artlSeq=5169402 file=extracted/1008/5169402/제27회 감정평가사 1차 문제/제27회 감정평가사 1교시문제(A형).pdf
%% conversion: kordoc 4.24.51 -> tools/gen1.py (auto); review=proofread-sonnet; flags=none
\begin{question}{18}
법률행위의 무효에 관한 설명으로 옳지 않은 것은? (다툼이 있으면 판례에 따름)
\begin{choices}
  \item 법률행위의 일부분이 무효인 경우, 그 무효부분이 없더라도 법률행위를 하였을 것이라고 인정될 때에는 나머지 부분은 무효가 되지 않는다.
  \item 매매계약이 매매대금의 과다로 인하여 불공정한 법률행위로서 무효인 경우, 무효행위의 전환에 관한 규정이 적용될 수 없다.
  \item 무효행위의 추인은 명시적으로뿐만 아니라 묵시적으로도 할 수 있다.
  \item 부동산 이중매매에서 매도인의 배임행위에 제2매수인이 적극 가담한 경우, 제2매수인의 매매계약은 무효이고 추인에 의하여 유효로 되지 않는다.
  \item 무효인 가등기를 유효한 등기로 전용하기로 한 약정은 그 때부터 유효하고, 이로써 그 가등기가 소급하여 유효한 등기로 전환될 수 없다.
\end{choices}
\end{question}
\answer{2}
